\documentclass[10pt,a4paper]{amsart}
\usepackage[margin=19mm]{geometry}
\usepackage{amsmath,amssymb,mathtools}
\usepackage[scr=rsfs,cal=euler]{mathalfa}
\usepackage{xfrac}
\usepackage[unicode,hidelinks,bookmarks=false]{hyperref}
\newcommand{\ie}{i.e.\ }
\newcommand{\cf}{cf.\ }
\newcommand{\resp}{respectively\ }
\newcommand{\Subsection}{\subsection}
\providecommand{\nobreakdash}{\nobreak\hbox{-}}
\setlength{\emergencystretch}{2em}
\multlinegap0pt
\usepackage{bm}
\usepackage{amssymb}
\usepackage{caption}
\usepackage[shortlabels]{enumitem}
\usepackage[all]{xy}
\usepackage{float}
\newtheorem{lemma}{Lemma}
\newtheorem{theorem}{Theorem}
\newtheorem{corollary}{Corollary}
\newtheorem{proposition}{Proposition}
\newtheorem{enumalph}{Enumalph}
\DeclareMathOperator*{\Ob}{Ob}
\DeclareMathOperator{\Part}{Part}
\DeclareMathOperator{\sign}{sign}
\newcommand{\stirling}[2]{\genfrac{\{}{\}}{0pt}0{#1}{#2}}
\newcommand{\tstirling}[2]{\genfrac{\{}{\}}{0pt}1{#1}{#2}}
\title{Monomial bialgebras}
\author{Arkady Berenstein, Jacob Greenstein, Jian-Rong Li}
\date{Source: arXiv:2602.02342v2 (CC BY 4.0)}
\begin{document}
\maketitle

\noindent\emph{Source and licence.} Monomial bialgebras. Version arXiv:2602.02342v2 (CC BY 4.0), distributed by arXiv under the Creative Commons Attribution 4.0 International licence, http://creativecommons.org/licenses/by/4.0/, as recorded in the arXiv licence metadata for this exact version and shown on the abstract page https://arxiv.org/abs/2602.02342v2. Full licence notice: \texttt{licenses/arxiv-2602.02342-CC-BY-4.0.txt}.\par\smallskip

\noindent\textit{Editorial scope.} Counted unit: the article's proposition trans (source0.tex lines 2668--2676). The article's complete original proof of it is delivered with it (source0.tex lines 2678--2856, one \texttt{proof} environment of 179 lines). No mathematics is written, repaired or re-derived: the statement and the proof are the article's own lines, in the article's own notation, and no retained line is substituted or re-flowed.  Retained uncounted as the source-local context the proof needs: lines 207--210; lines 2417--2428; lines 2542--2547; lines 2584--2586.  Nothing was omitted from any retained span, so the package carries no omission record.  No retained line contains a citation, so the wrapper carries no bibliography and no bibliography entry.  No retained line contains a figure environment, an image-inclusion command or drawing code, so no image is delivered and the package declares no asset dependency.  Editorial formatting only, in addition: an AMS \texttt{amsart} preamble that restates the article's own declarations and macros used by the retained lines, namely the environments \texttt{lemma}, \texttt{theorem}, \texttt{corollary}, \texttt{proposition}, \texttt{enumalph} on separate counters, together with the standard packages those lines need.  The retained environments print in this document as Lemma 1, Theorem 1, Corollary 1, Proposition 1, Enumalph 1 in order of appearance, whereas the article prints the same items with the numbering of its own full document.  The complete native source of the article as distributed in the arXiv e-print for this exact version is bundled unchanged as \texttt{sources/193845/source0.tex}.
\begin{equation}\label{eq:eps w alpha defn}
\boldsymbol\epsilon(w,\boldsymbol d)_{i,j}=\delta_{i,j}d_i+\sign(w(j)-w(i)),
\qquad 1\le i,j\le n.
\end{equation}

\begin{lemma}\label{lem:In trans to trans}
Let~$w\in S_n$, $\boldsymbol{\alpha}:[n]\to \{1,-1\}$ and define $\boldsymbol\epsilon(w,\boldsymbol{\alpha}):[n]\times[n]\to \{1,-1\}$ by
$$
\boldsymbol\epsilon(w,\boldsymbol{\alpha})(i,j)=\delta_{i,j}\alpha(i)+
\sign(w(j)-w(i)),\qquad i,j\in [n]
$$
(cf.~\eqref{eq:eps w alpha defn}).
Then~$\boldsymbol\epsilon(w,\boldsymbol{\alpha})$ is transitive. Moreover,
all transitive maps~$\mathbf c:[n]\times[n]\to\{1,-1\}$ satisfying
$\mathbf c(j,i)=-\mathbf c(i,j)$, $(i,j)\in I_n$, are obtained this way.
In particular, there are $n! 2^n$ such maps.
\end{lemma}

\begin{theorem}
\label{th:polynomial} 
Let~$M$ be a finite set and $F$ be an $M$-admissible functor $\mathbf{Set}_0\to \mathbf{Set}$. Then
$|F(C)|=p_F(|C|)$ 
for any finite $C\in \Ob(\mathbf{Set}_0)$.
\end{theorem}

\begin{corollary}\label{cor:stirling F2}
$\tstirling{M}{2}_F=\frac{1}{2}|F([2])|-1$ under the assumptions of Theorem \ref{th:polynomial}. 
\end{corollary}

\begin{proposition}\label{prop:enum trans}
For~$\Gamma=\overset{\leftrightarrow\circlearrowleft}{K_n}$, we have $$
\stirling{\Gamma}{k}=\begin{cases}2\stirling{n}{k},& k\ge 3,\\\tfrac12 B_n-1,&k=2,
\end{cases}
$$
where
$B_n :=\sum\limits_{1\le k\le n} \tstirling{n}{k}2^k k!$
is the number of transitive relations~$\mathscr R$ on~$[n]$ such that~$([n]\times [n])\setminus\mathscr R$ is also transitive.
\end{proposition}

\begin{proof}
Given~$i,j\in[n]$ and~$\mathcal P\in\Part([n]\times[n])$,
abbreviate~$\mathcal P(i,j):=\mathcal P[(i,j)]$.
Let~$\mathscr T_n$ be the set of all $\mathcal P\in \Part([n]\times[n])$
such that~$\mathcal P(i,k)\in\{\mathcal P(i,j),\mathcal P(j,k)\}$ for all~$i,j,k\in[n]$. 
It is immediate from the definition that~$\Part_r^{\mathcal E_\Gamma([r])}([n]\times[n])
=\mathscr T_n\cap \Part_r([n]\times[n]):=\mathscr T_{n,r}$.


Given~$\mathcal P\in\Part([n]\times[n])$, define a relation~$\sim_{\mathcal P}$
on~$[n]$ by~$i\sim_{\mathcal P}j$ if and only if $\mathcal P(i,j)=\mathcal P(j,i)$,
$i,j\in[n]$. This relation is clearly reflexive and symmetric.
\begin{lemma}\label{lem:prop sim_P}
Let~$\mathcal P\in\mathscr T_n$. Then
\begin{enumalph}
\item\label{lem:prop sim_P.a} $\sim_{\mathcal P}$ is an equivalence relation;
\item\label{lem:prop sim_P.c} If~$i\sim_{\mathcal P}i'$,
$j\sim_{\mathcal P}j'$ then~$\mathcal P(i,j)=\mathcal P(i',j')$.
\end{enumalph}
\end{lemma}
\begin{proof}
Note that if~$\mathcal P\in\mathscr T_n$ and~$\mathcal P(i,j)=\mathcal P(j,i)$
then~$\mathcal P(i,i)=\mathcal P(i,j)=\mathcal P(j,i)=\mathcal P(j,j)$.
Suppose that~$i\sim_{\mathcal P}j$, $j\sim_{\mathcal P}k$, $i,j,k\in[n]$. Then~$\mathcal P(i,i)=\mathcal P(j,j)=\mathcal P(k,k)$,
$\mathcal P(i,k)\in\{\mathcal P(i,j),\mathcal P(j,k)\}=\{\mathcal P(j,j)\}$
and~$\mathcal P(k,i)\in\{\mathcal P(k,j),\mathcal P(j,i)\}=\{\mathcal P(j,j)\}$.
Therefore $\mathcal P(i,k)=\mathcal P(k,i)$, that is~$i\sim_{\mathcal P}k$.
To prove~\ref{lem:prop sim_P.c}, 
suppose that~$\mathcal P(i,j)\not=\mathcal P(i',j')$. Since
$\mathcal P(i,j)\in\{\mathcal P(i,i'),\mathcal P(i',j)\}$ while
$\mathcal P(i',j')\in\{\mathcal P(i',i),\mathcal P(i,j')\}=
\{\mathcal P(i,i'),\mathcal P(i,j')\}$, it follows that~$\mathcal P(i,j)=
\mathcal P(i',j)\not=\mathcal P(i,j')=\mathcal P(i',j')$. Yet,
$\mathcal P(i',j)\in\{\mathcal P(i',j'),\mathcal P(j',j)\}$, whence~$\mathcal P(i',j)=\mathcal P(j,j')$, while~$\mathcal P(i,j')\in\{\mathcal P(i,j),\mathcal P(j,j')\}$ whence~$\mathcal P(i,j')=\mathcal P(j,j')$. Thus, $\mathcal P(i,j')=\mathcal P(i',j)$ which is a contradiction.
\end{proof}
We need the following
\begin{lemma}\label{lem:surj map part}
Let~$\pi:[n]\to[r]$
be surjective. 
\begin{enumalph}
    \item\label{lem:surj map part.a} Suppose that $\mathcal P\in\Part([n]\times[n])$
    satisfies 
    $\mathcal P(i,j)=\mathcal P(i',j')$ for all $i,i',j,j'\in[n]$ such that $\pi(i)=\pi(i')$ and~$\pi(j)=\pi(j')$.
Then~$(\pi\times\pi)(\mathcal P):=\{ (\pi\times\pi)(P)\,:\, P\in\mathcal P\}\in \Part([r]\times[r])$;
\item\label{lem:surj map part.b} Suppose that~$\mathcal Q\in \Part([r]\times[r])$.
Then~$(\pi\times\pi)^{-1}(\mathcal Q):=\{ (\pi\times\pi)^{-1}(Q)\,:\, Q\in\mathcal Q\}$
is a partition of~$[n]\times[n]$ and~$(\pi\times\pi)((\pi\times\pi)^{-1}(\mathcal Q))=\mathcal Q$.
\end{enumalph}
\end{lemma}
\begin{proof}
Since~$\bigcup_{P\in\mathcal P} P=[n]\times[n]$ and~$\pi$
is surjective, $\bigcup_{P\in\mathcal P}
(\pi\times\pi)(P)=[r]\times[r]$. Suppose that~$(i,j)\in 
(\pi\times\pi)(P)\cap(\pi\times\pi)(P')$, $i,j\in[r]$. Then~$i=(\pi(i'),
\pi(j'))=(\pi(i''),\pi(j''))$ for some~$i',i'',j',j''\in[n]$
and~$P=\mathcal P(i',j')=\mathcal P(i'',j'')=P'$
whence~$(\pi\times\pi)(P)=(\pi\times\pi)(P')$.
Thus, $(\pi\times\pi)(\mathcal P)\in\Part([r]\times[r])$.

Part~\ref{lem:surj map part.b} is immediate.
\end{proof}


Given~$\mathcal P\in\mathscr T_n$, let~$\mathcal C(\mathcal P)$ be
the partition of~$[n]$ into equivalence classes for~$\sim_{\mathcal P}$. Denote~$\underline{\mathscr T}_n$ the set of 
all~$\mathcal P\in\mathscr T_n$ satisfying~$\mathcal P(i,j)\not=\mathcal P(j,i)$ for all~$i\not=j\in[n]$.
\begin{lemma}\label{lem:surj trans part}
Let~$\mathcal P\in\mathscr T_n$ and suppose that~$\mathcal C(\mathcal P)\in\Part_r([n])$. Then
$(\pi\times\pi)(\mathcal P)\in \underline{\mathscr T}_r$ 
for any surjective~$\pi:[n]\to[r]$ such that~$\mathcal C(\mathcal P)=\{ \pi^{-1}(i)\}_{i\in[r]}$.
\end{lemma}
\begin{proof}
By
Lemmata~\ref{lem:prop sim_P}\ref{lem:prop sim_P.c} and~\ref{lem:surj map part}, 
$(\pi\times\pi)(\mathcal P)\in\Part([r]\times[r])$. Let~$i=\pi(i'),j=\pi(j'),k=\pi(k')\in[r]$, $i',j',k'\in[n]$.
Then~$\mathcal P(i',k')\in \{ \mathcal P(i',j'),\mathcal P(j',k')\}$,
and so $(\pi\times\pi)(\mathcal P)(i,k)\in\{(\pi\times\pi)(\mathcal P)(i,j),
(\pi\times\pi)(\mathcal P)(j,k)\}$, that is~$(\pi\times\pi)(\mathcal P)\in\mathscr T_r$.
Finally, suppose that~$(\pi\times\pi)(\mathcal P)(i,j)=(\pi\times\pi)(\mathcal P)(j,i)$ 
for some~$i\not=j\in[r]$. Then
$i=\pi(i')$, $j=\pi(j')$, $i',j'\in[n]$ with~$i'\not\sim_{\mathcal P}j'$
yet~$(\pi\times\pi)(\mathcal P(i',j'))=(\pi\times\pi)(\mathcal P(j',i'))$, that
is, $(i,j)=(\pi(i''),\pi(j''))$ for some~$i'',j''\in[n]$ such that~$\mathcal P(i'',j'')=
\mathcal P(j',i')\not=\mathcal P(i',j')$. Since~$\mathcal P(i'',j'')=\mathcal P(i',j')$
by Lemma~\ref{lem:prop sim_P}\ref{lem:prop sim_P.c}, this is a contradiction.
\end{proof}

Let~$\mathcal J\in\Part_r([n])$.
Define~$\mathcal J_k=\mathcal J[m_k]$ where~$m_k=\min\big([n]
\setminus \bigcup_{1\le t<k} \mathcal J_t\big)$; thus, $m_1=1$
and~$\mathcal J_1=\mathcal J[1]$ and so on.
Let~$\pi_{\mathcal J}$ be the unique surjective map~$\pi:[n]\to[r]$
satisfying~$\pi^{-1}(t)=\mathcal J_t$, $t\in[r]$. 
By Lemma~\ref{lem:surj trans part} 
we obtain a well-defined map~$\Pi_n:\mathscr T_n\to \bigsqcup_{1\le r\le n}
\underline{\mathscr T}_r$, $\mathcal P\mapsto (\pi_{\mathcal C(\mathcal P)}\times \pi_{\mathcal C(\mathcal P)})(\mathcal P)$,
$\mathcal P\in\mathscr T_n$.
\begin{lemma}\label{lem:collapse part}
The map~$\Pi_n$ is surjective and~$|\Pi_n^{-1}(\mathcal P)|=\tstirling{n}{r}$ 
for any~$\mathcal P\in\underline{\mathscr T}_r$. 
\end{lemma}
\begin{proof}
Let~$\mathcal P\in\underline{\mathscr T}_r$, 
$\mathcal J\in\Part_r([n])$ and define~$\mathcal Q:=\mathcal Q(\mathcal P,\mathcal J)=(\pi_{\mathcal J}\times\pi_{\mathcal J})^{-1}(\mathcal P)\in \Part([n]\times[n])$ by Lemma~\ref{lem:surj map part}\ref{lem:surj map part.b} and $(\pi_{\mathcal J}\times\pi_{\mathcal J})(\mathcal Q)=\mathcal P$.
We claim that~$\mathcal Q\in \mathscr T_n$. Indeed,
let~$i,j,k\in[n]$. Then
$$
(\pi_{\mathcal J}\times\pi_{\mathcal J})(\mathcal Q(i,k))=\mathcal P(\pi_{\mathcal J}(i),\pi_{\mathcal J}(k))\in 
\{ \mathcal P(\pi_{\mathcal J}(i),\pi_{\mathcal J}(j)),\mathcal P(\pi_{\mathcal J}(j),\pi_{\mathcal J}(k))\},
$$
whence~$\mathcal Q(i,k)\in\{ \mathcal Q(i,j),
\mathcal Q(j,k)\}$. 

Next, we prove that~$\mathcal C(\mathcal Q)=\mathcal J$,
that is, $\mathcal Q(i,j)=\mathcal Q(j,i)$ if and only if~$\pi_{\mathcal J}(i)=\pi_{\mathcal J}(j)$. Suppose first that~$\mathcal Q(i,j)=\mathcal Q(j,i)$. Then $\mathcal P(\pi_{\mathcal J}(i),\pi_{\mathcal J}(j))=
\mathcal P(\pi_{\mathcal J}(j),\pi_{\mathcal J}(i))
$ which forces~$\pi_{\mathcal J}(i)=\pi_{\mathcal J}(j)$
since~$\mathcal P\in\bigsqcup_{r\in[n]}\underline{\mathscr T}_r$. 
The converse is obvious
since~$\mathcal Q(i,j)=(\pi_{\mathcal J}\times\pi_{\mathcal J})^{-1}(
\mathcal P(\pi_{\mathcal J}(i),\pi_{\mathcal J}(j))$, $i,j\in[n]$.

Thus, we proved that~$\Pi_n(\mathcal Q(\mathcal P,\mathcal J))=
\mathcal P$, that is, $\Pi_n$ is surjective and, moreover,
that~$\Pi^{-1}_n(\mathcal P)=\{ \mathcal Q(\mathcal P,\mathcal J)\,:\,
\mathcal J\in \Part_r([n])\}$. It remains to observe that,
since $\mathcal J=\mathcal C(\mathcal Q(\mathcal P,\mathcal J))$,
the map $\Part_r([n])\to \Pi_n^{-1}(\mathcal P)$, $\mathcal J\mapsto 
\mathcal Q(\mathcal P,\mathcal J)$, $\mathcal J\in\Part_r([n])$
is also injective.
\end{proof}
\begin{lemma}\label{lem:trans count last}
Let~$\mathcal P\in \underline{\mathscr T}_n$.
\begin{enumalph}
\item\label{lem:trans count last.a} If~$\mathcal P(i,i)\not=\mathcal P(j,j)$ for all~$i\not=j\in[n]$
then $\mathcal P=\{ \mathcal P_1,\dots,\mathcal P_n\}$ where 
either~$\mathcal P_i=\{ (i,j)\,:\, j\in[n]\}$ for all~$i\in [n]$
or~$\mathcal P_i=\{ (j,i)\,:\, j\in [n]\}$ for all~$i\in[n]$.

\item\label{lem:trans count last.b} If~$\mathcal P(i,i)=\mathcal P(j,j)$ for some~$i\not=j$ then~$\mathcal P\in \Part_2([n])$. 
\end{enumalph}
\end{lemma}
\begin{proof}
To prove~\ref{lem:trans count last.a}, 
note that for~$\mathcal P\in\underline{\mathscr T}_n$
and satisfying~$\mathcal P(i,i)\not=\mathcal P(j,j)$ for all $j\not=i\in[n]$, since
$\{\mathcal P(i,i),\mathcal P(j,j)\}\subset \{\mathcal P(i,j),\mathcal P(j,i)\}$
it follows that, without loss of generality, 
$\mathcal P(i,i)=\mathcal P(i,j)$ and~$\mathcal P(j,j)=\mathcal P(j,i)$.
If~$n=2$ we are done. If~$n>2$, let~$k\in[n]\setminus\{i,j\}$
and suppose first that $\mathcal P(i,k)=\mathcal P(k,k)$, whence $\mathcal P(k,i)=\mathcal P(i,i)$. Then $\mathcal P(i,i)=\mathcal P(k,i)\in
\{\mathcal P(k,j),\mathcal P(j,i)\}=\{\mathcal P(k,j),\mathcal P(j,j)\}
\subset \{\mathcal P(j,j),\mathcal P(k,k)\}$ which is a contradiction. Thus,
$\mathcal P(i,k)=\mathcal P(i,i)$ and then $\mathcal P(k,i)=\mathcal P(k,k)$. As $\mathcal P(k,k)=\mathcal P(k,i)\in\{\mathcal P(k,j),\mathcal P(j,i)\}
=\{\mathcal P(k,j),\mathcal P(j,j)\}$ and $\mathcal P(k,j)\in\{\mathcal P(j,j),\mathcal P(k,k)\}$ this forces $\mathcal P(k,j)=\mathcal P(k,k)$ and $\mathcal P(j,k)=\mathcal P(j,j)$.
Therefore, $\mathcal P(i,i)=\mathcal P(i,j)=\mathcal P(i,k)$, $\mathcal P(j,i)=\mathcal P(j,j)=\mathcal P(j,k)$ and $\mathcal P(k,i)=\mathcal P(k,j)=\mathcal P(k,k)$.
This implies that $\mathcal P(r,s)=\mathcal P(r,r)$ for all $r,s\in[n]$. 

To prove part~\ref{lem:trans count last.b}, suppose that
$\mathcal P(i,i)=\mathcal P(j,j)$ for some $i\not=j$.
Since $\mathcal P(i,i)\in\{\mathcal P(i,j),\mathcal P(j,i)\}$
and $\mathcal P(i,j)\not=\mathcal P(j,i)$, we may assume without loss of generality that $\mathcal P(i,j)=\mathcal P(i,i)=\mathcal P(j,j)\not=\mathcal P(j,i)$.

Let $k\not=i$. 
Suppose that $\mathcal P(i,k)\notin\{\mathcal P(i,j),\mathcal P(j,i)\}$. Then $\mathcal P(i,k)\in\{\mathcal P(i,j),\mathcal P(j,k)\}$
forces $\mathcal P(j,k)=\mathcal P(i,k)$ while $\mathcal P(i,j)=\mathcal P(i,i)\in\{\mathcal P(i,k),\mathcal P(k,i)\}$ yields $\mathcal P(k,i)=\mathcal P(i,j)$.
But then $\mathcal P(j,i)\in\{\mathcal P(j,k),\mathcal P(k,i)\}=\{\mathcal P(i,j),
\mathcal P(i,k)\}$, which is a contradiction.
If $\mathcal P(k,i)\notin\{\mathcal P(i,j),\mathcal P(j,i)\}$, then
$\mathcal P(j,i)\in\{\mathcal P(j,k),\mathcal P(k,i)\}$ forces $\mathcal P(j,k)=\mathcal P(j,i)$.
Since $\mathcal P(i,j)=\mathcal P(j,j)\in\{\mathcal P(j,k),\mathcal P(k,j)\}$ we have $\mathcal P(k,j)=\mathcal P(i,j)$. Then $\mathcal P(k,i)\in\{\mathcal P(k,j),\mathcal P(j,i)\}
=\{\mathcal P(i,j),\mathcal P(j,i)\}$ which is a contradiction. 

Therefore, for any~$k\in[n]\setminus\{i,j\}$, $\mathcal P(i,k),\mathcal P(k,i)\in \{\mathcal P(i,j),\mathcal P(j,i)\}$.
But then $\mathcal P(k,l)\in\{\mathcal P(k,i),\mathcal P(i,l)\}\subset \{\mathcal P(i,j),\mathcal P(j,i)\}$ for all $k,l\in [n]$. Thus, $\mathcal P\in\Part_2([n]\times[n])$. 
\end{proof}
It follows from Corollary~\ref{cor:stirling F2} and Lemmata~\ref{lem:collapse part}, \ref{lem:trans count last}\ref{lem:trans count last.a} and~\ref{lem:In trans to trans} that~$|\mathscr T_{n,r}|=2\tstirling{n}{r}$,
while~$\mathcal E_{\Gamma}([2])=\sum_{1\le r\le n}\tstirling{n}{r}2^r r!=B_n$. The assertion is now immediate.
\end{proof}

\end{document}
